\section{Introduction}

Let $p_n$ denote the $n$th prime and write $d_n=p_{n+1}-p_n$.
The prime number theorem identifies $\log p$ as the average scale of
prime gaps near $p$. Individual gaps can be much larger.
Westzynthius proved that $d_n/\log p_n$ is unbounded
\cite{Westzynthius1931}; the constructions of Erd\H{o}s and Rankin
gave increasingly strong quantitative lower bounds for exceptionally
large gaps \cite{Erdos1935,Rankin1938}.
Modern refinements were obtained independently by Ford, Green,
Konyagin and Tao and by Maynard
\cite{FordGreenKonyaginTao2016,Maynard2016Large}, followed by their
joint work \cite{FordGreenKonyaginMaynardTao2018}.
These results concern the size of very large gaps. A lower bound for
the largest gap does not by itself give a positive proportion of
gaps exceeding a fixed multiple of $\log p$.

We study this frequency question with equal weight on the prime
indices. For a set $A$ of positive integers, its lower asymptotic
density is $\liminf_{N\to\infty}|A\cap[1,N]|/N$.

\begin{theorem}\label{thm:main}
For every fixed real $C>0$, there are constants $c(C)>0$ and $N_0(C)$
such that, for every integer $N\ge N_0(C)$,
\[
 \#\{1\le n\le N:p_{n+1}-p_n>C\log p_n\}\ge c(C)N.
\]
\end{theorem}

The density is ordinary lower density: the bound holds for every
sufficiently large initial segment, and both constants may depend
on the fixed threshold $C$.

One consequence concerns successive values of $p_n/n$.
Erd\H{o}s and Prachar asked whether the indices at which this ratio
increases have positive lower density \cite[p.~256]{ErdosPrachar1962}.

\begin{corollary}\label{cor:ratios}
The set
\[
 \left\{n\ge1:\frac{p_n}{n}<\frac{p_{n+1}}{n+1}\right\}
\]
has positive lower asymptotic density.
\end{corollary}
\begin{proof}
The displayed inequality is equivalent to
$p_{n+1}-p_n>p_n/n$.
The prime number theorem gives $p_n/n\sim\log p_n$.
Apply Theorem~\ref{thm:main} with $C=2$ and discard the finitely many
indices for which $p_n/n\ge2\log p_n$.
\end{proof}

Thus the corollary answers the Erd\H{o}s--Prachar question affirmatively.

\subsection*{Gap frequency and empty intervals}

The distinction between counting intervals and counting prime gaps is
essential. A long gap contains many integer starting points of
prime-free short intervals. Consequently a positive proportion of
empty intervals may be accounted for by a sparse set of very long
gaps. Bazzanella, Languasco and Zaccagnini distinguish the two counting
measures explicitly \cite{BLZ2010}. Their Theorem~5 gives positive
prime-start proportions for each fixed threshold below
$2/3.454=0.579038\ldots$, whereas their Theorem~3 reaches some
intervals longer than the average prime spacing by counting integer
starts. Tao's empty-interval argument \cite{Tao2019} also identifies
the additional control required to pass from interval measure to
gap counts.

Conditional results describe much more of the distribution.
Gallagher derived Poisson statistics for prime counts in
integer-start intervals of length $\lambda\log X$, for fixed
$\lambda>0$, from a uniform
Hardy--Littlewood prime-tuple hypothesis
\cite[Theorem~1]{Gallagher1976}.
The Poisson-tail results of Jha
\cite[Theorems~1.4 and~1.6]{Jha2026} treat growing parameter ranges
under a stronger uniform tuple hypothesis.
The counting measures and uniformity assumptions matter in comparing
these conclusions with Theorem~\ref{thm:main}.

\subsection*{The adjacent-interval construction}

We control the number of distinct gaps by requiring a prime just
before an empty interval. Fix $\lambda>0$, put
$h=\lfloor\lambda\log X\rfloor$, and average over integers
$X<m\le2X$. We construct a nonnegative weight for which a prime in
$(m,m+h]$ has substantial weighted probability, whereas the expected
number of primes in $(m+h,m+2h]$ is arbitrarily small.
One second-moment estimate detects the first prime; a separate
second-moment estimate converts the resulting weighted mass into a
positive proportion of ordinary integer starts.
The last prime in $(m,m+h]$ then begins a gap longer than $h$.
It can be selected by at most $h$ values of $m$, independently of
the length of that gap. This bounded multiplicity is the step that
turns the interval construction into a count of consecutive gaps.

The weights are squares of signed sums of smooth divisor sums, each
built from divisors of shifted integers in the second interval.
The test functions make terms involving different shift sets
orthogonal in the limiting unmarked second moment.
At a prime belonging to one of those shift sets, the support cutoff
forces the corresponding divisor coordinate to equal one in every
nonzero summand. Removing that shift then couples terms whose
sets differ in size by one. The central choice is an alternating
family for which these neighboring terms nearly cancel.
The resulting square has positive average but arbitrarily small
weighted prime mass in the second interval. A detector estimate
for the first interval has a constant independent of the size of
the chosen family, which allows the suppression to be made strong
enough before the final counting argument.

The analytic ingredients are smooth divisor-sum correlations, the
Bombieri--Vinogradov distribution theorem, and an average of the
prime-tuple singular series. The distribution theorem originates
in the work of Bombieri and A.~I.~Vinogradov
\cite{Bombieri1965,Vinogradov1965}.
Goldston--Y{\i}ld{\i}r{\i}m's correlation method is a direct
antecedent: they square a linear combination of tuple approximations
of different cardinalities
\cite[p.~6, equation~(2.14) and the preceding paragraph]{GY2005}.
Maynard's smooth multidimensional sieve and product-profile
construction \cite{Maynard2015} provide further antecedents.
Here the orthogonality and coordinate-removal identities reduce the
choice of a weight to cancellation between adjacent dimensions.
All dimensions and test functions are fixed before $X$ tends to
infinity; the uniform error estimates needed to sum over the shifts
are proved in the paper.

Section~\ref{sec:counting} first deduces the theorem from a precise
weight proposition and proves the passage to ordinary index density.
Section~\ref{sec:weights} defines the square and states its moment
identities; Section~\ref{sec:cancellation} chooses the test functions
that make the adjacent dimensions cancel.
Section~\ref{sec:moments} proves the general mixed-moment estimates,
the singular-series average, and the moment identities used in the
construction.
